\chapter{Current v0.0.7 Evaluation and Validation}
\label{chap:v007}

\section{Evidence sets and denominators}
This chapter updates the historical v0.0.6 evaluation in
Chapter~\ref{chap:results}. Three v0.0.7 evidence sets answer different
questions and must retain their own denominators. The production diagnostic
scan froze 207 repositories at recorded commits; 206 yielded analyzable
outputs and one timed out. A paired semantic change was evaluated on those
same 206 outputs. The separate numerical-threshold benchmark parsed all 207
frozen repositories and extracted unique metric observations. Its zero scan
failures do not retroactively make the production provider run complete.

\begin{table}[htbp]
\centering\small
\caption{Scope of the current evidence sets. Counts are observations, not validated defects.}
\label{tab:v007_scope}
\begin{tabularx}{\textwidth}{p{3.1cm}p{3.0cm}X}
\toprule
Evidence set & Population & What it establishes \\
\midrule
Production scan & 206 analyzable / 207 frozen repositories & Source/provider
counts and execution health; one repository timeout persists. \\
Semantic paired run & Same 206 analyzable repositories & Effect of required
role-argument reasoning with fixed commits, file manifests, profile and
provider scope. \\
Threshold benchmark & 207 / 207 repositories & Distributions and candidate
rates for 24 numerical metrics; no human accuracy estimate yet. \\
Agent-labelled review & 847 sampled units & Diagnostic calibration of
variable resolution; one reviewer, with unresolved cases retained. \\
\bottomrule
\end{tabularx}
\end{table}

\section{Production execution and rule composition}
The v0.0.7 diagnostic run produced 390,441 raw findings and 376,271 canonical
occurrences. Of the 207 frozen repositories, 171 had complete required and
optional provider output, 35 were partial because of ansible-lint, and one
failed after a 1,800-second timeout. The same repository timed out on the
paired candidate run even with a 3,600-second allowance. Thus the 376,271
total describes the 206 analyzable repositories; the missing repository is
never encoded as zero findings. GitLeaks completed for all 206 analyzable
repositories. The 108,757 ansible-lint occurrences across all outputs include
partial runs; the provider-comparable subset contains 83,419 occurrences on
the 171 complete repositories.

\begin{figure}[htbp]
\centering
\begin{tikzpicture}[x=1cm,y=1cm]
  \fill[anspaper,rounded corners=4pt] (0,0) rectangle (13.3,4.25);
  \node[anchor=west,text=ansnavy,font=\bfseries\large] at (.45,3.75)
    {207 frozen repositories};
  \node[anchor=west,text=ansgray,font=\small] at (.45,3.28)
    {Production v0.0.7 execution state};
  \fill[ansblue!80] (.45,2.45) rectangle (11.04,2.94);
  \fill[ansgold!88] (11.04,2.45) rectangle (13.20,2.94);
  \fill[ansred!78] (13.20,2.45) rectangle (13.26,2.94);
  \node[anchor=west,text=white,font=\bfseries] at (.65,2.70) {171 complete};
  \node[anchor=east,text=ansnavy,font=\bfseries\small] at (13.08,2.70) {35 partial};
  \draw[ansred,thick] (13.23,2.35)--(12.53,1.82);
  \node[anchor=east,text=ansred,font=\small\bfseries] at (12.48,1.82)
    {1 timed out};
  \node[anchor=west,font=\small,text=ansnavy] at (.45,1.32)
    {206 analyzable outputs};
  \node[anchor=west,font=\small,text=ansgray] at (.45,.83)
    {Partial provider evidence stays visible and is excluded from comparable estimates.};
\end{tikzpicture}
\caption{Repository health in the v0.0.7 production scan. The failed
repository is missing data, not an observed zero.}
\label{fig:v007_health}
\end{figure}

The internal selected-source population contributes 265,586 occurrences in
the study report. Three rules account for 174,847, or 65.8\% of that internal
population: \code{missing-tags} (98,005), \code{unguarded-variable}
(52,721), and \code{variable-naming-convention} (24,121). The first is an
optional, low-confidence policy signal. A large count from a permissive rule
does not imply the same number of confirmed defects. The Git-history review
population contributes 494 occurrences under a different unit and denominator;
it is excluded from selected-source rates. A single repository contributes
2,776 \code{yaml[quoted-strings]} ansible-lint findings, illustrating why
pooled totals must be accompanied by repository prevalence.

\section{Paired role-contract correction}
The candidate semantic run adds explicit handling of required role arguments.
It reduced canonical occurrences from 376,271 to 376,024 and uncertain
occurrences from 54,443 to 54,196. Ansible-lint and GitLeaks totals remained
unchanged. Comparison identities initially marked 568 fixed and 321
introduced; 320 of the internal ``introduced'' records share repository,
file, and line with a fixed record and represent a change from
\code{unresolved-runtime-input} to \code{undeclared-input-contract}. After
pairing those reclassifications, the net detector effect is exactly 247
suppressed \code{unguarded-variable} occurrences. One further fixed/introduced
pair belongs to an ansible-lint \code{warning[outdated-tag]} identity change.

\begin{figure}[htbp]
\centering
\begin{tikzpicture}[x=1cm,y=1cm]
  \fill[anspaper,rounded corners=4pt] (0,0) rectangle (13.3,4.3);
  \node[anchor=west,text=ansnavy,font=\bfseries\large] at (.45,3.76)
    {What changed in variable findings};
  \node[draw=ansblue,fill=white,rounded corners=3pt,minimum width=3cm,
    minimum height=1.15cm,align=center] (old) at (2.2,2.33)
    {567 old identities\\marked fixed};
  \node[draw=ansteal,fill=white,rounded corners=3pt,minimum width=3cm,
    minimum height=1.15cm,align=center] (reclass) at (7.7,2.95)
    {320 reclassified\\same source lines};
  \node[draw=ansgold,fill=white,rounded corners=3pt,minimum width=3cm,
    minimum height=1.15cm,align=center] (removed) at (7.7,1.58)
    {247 suppressed\\required inputs};
  \draw[-{Latex[length=2.2mm]},thick,ansnavy] (old.east)--(reclass.west);
  \draw[-{Latex[length=2.2mm]},thick,ansnavy] (old.east)--(removed.west);
  \node[anchor=west,font=\small,text=ansgray] at (.48,.72)
    {Net canonical change: 376,271 $\rightarrow$ 376,024 on the same 206 analyzable repositories.};
\end{tikzpicture}
\caption{Reconciliation of comparison identities with source-level semantic
changes. The 320 relabelled cases are existing findings.}
\label{fig:v007_variable_change}
\end{figure}

The 247 suppressions occur in seven repositories, including 121 in
Linuxfabrik/lfops, 44 in redhatci/ansible-collection-redhatci-ocp, 33 in
Dell's OpenManage collection, and 32 in Alfresco's deployment repository.
Source inspection confirms the intended boundary: a \code{required: true}
role argument is treated as guaranteed at direct role entry; optional
arguments stay visible with an uncertainty reason. The candidate retains 997
\code{missing-internal-variable} occurrences, changes
\code{undeclared-input-contract} from 34,006 to 34,326, and changes
\code{unresolved-runtime-input} from 17,718 to 17,151. These are semantic
classes, not manual correctness labels.

\section{Variable-resolution review}
The frozen primary review has 847 units: a census of 247 removed required
argument findings, three 150-occurrence samples for candidate missing,
external-input uncertainty, and runtime-context uncertainty, plus 150
source-first files. The single agent reviewer verified 247 safe suppressions.
For \code{CANDIDATE_MISSING}, it labelled 84 true positives, 66 uncertain,
and zero false positives. The decided-case ratio is 84/84, whereas the
conservative true-positive yield across all 150 sampled candidates is 56.0\%.
The apparent 100\% decided-case result must not be presented as independent
human precision. Of 450 sampled classification cases, 366 received decided
labels, and all 366 agreed with the proposed resolution class.

Source-first review identified 11 supported issue occurrences and matched 10,
an observed 10/11 occurrence ratio. Those issues came from only two
repositories; one repository contained the unmatched case. Neither ratio is a
general recall estimate. The missed \code{tarzan_proj_dir} use in Miarka's
\code{roles/tarzan/defaults/main.yml}, beside a defined
\code{tarzan_root_dir}, is a useful false-negative case for detector repair.
The review used one agent and source context after sampling; independent
human agreement, adjudication, and unseen-repository validation remain open.

\section{Numerical threshold benchmark}
The corrected benchmark extracts 2,986,459 unique scalar observations for 24
metrics from all 207 frozen repositories, with zero scan failures and zero
duplicate observation identities. Detector-relevant denominators differ by
metric: 128,220 task observations come from 193 repositories, while 21,455
playbook-family observations come from 197. The benchmark repaired duplicate
manifest paths, excluded vars/metadata documents from playbook metrics,
replayed complete composite detector decisions, and retained parser source
positions for repeated task names. It generated 16,971 primary review rows
and 75-row agreement pilots for each metric (1,800 rows per reviewer).

Large-task size illustrates the method. Its distribution has Q1 = 4,
median = 6, Q3 = 10, and Tukey upper fence = 19 lines. The current cutoff of
15 is also the measured 90th percentile. Moving the cutoff to Q3 would
produce 27,158 predictions, versus 12,202 at the current cutoff: 14,956 more
review items, or 2.23 times the present workload. A threshold is therefore a
review-cost decision requiring human labels; a box-plot fence alone cannot
select it.

\begin{figure}[htbp]
\centering
\begin{tikzpicture}
\begin{axis}[
  xbar,width=.89\textwidth,height=7.2cm,
  xmin=0,xmax=32000,bar width=12pt,
  symbolic y coords={P95,Tukey,current,P80,Q3},ytick=data,
  yticklabel style={font=\small},
  xlabel={Predicted tasks among 128,220 eligible observations},
  xmajorgrids=true,grid style={ansgray!20},
  axis line style={ansnavy!50},tick style={ansnavy!50},
  scaled x ticks=false,
  xtick={0,8000,16000,24000,32000},
  xticklabel style={/pgf/number format/fixed,/pgf/number format/1000 sep={,}},
  nodes near coords,nodes near coords style={font=\scriptsize,anchor=west},
  enlarge y limits=0.16]
\addplot[fill=ansblue!75,draw=none] coordinates
  {(27158,Q3) (22566,P80) (12202,current) (7738,Tukey) (6378,P95)};
\end{axis}
\end{tikzpicture}
\caption{Large-task candidate workload. Q3 = 10 LOC, P80 = 11,
current/P90 = 15, Tukey fence = 19, and P95 = 21. Candidate quality is
unmeasured until independent labels are completed.}
\label{fig:v007_threshold_sensitivity}
\end{figure}

The current production settings produce 11,124 complex-expression decisions,
12,202 large-task decisions, 3,104 active large-playbook decisions, 1,451
large-play decisions, and 52 special-character naming decisions in the
corresponding benchmark populations. These are separate metric studies and
must not be summed. The accepted \code{avg_task_playbook} cutoff is measured
but does not affect the production large-playbook decision;
\code{similarityPlays} is an accepted but unused configuration key. This
inventory prevents a purported calibration of a threshold with no production
effect.

\section{Validation gate and interpretation}
The threshold packets are prepared for two independent human reviewers.
The agreement pilot records independent positive, negative, or uncertain
labels; Cohen's kappa (target 0.80), Gwet's AC1, prevalence, and label-specific
agreement are calculated before codebook freeze. The primary stratified
sample carries inclusion probabilities and inverse-probability weights.
After adjudication, candidates can be scored with weighted precision, recall,
F1 and F0.5 and project-cluster bootstrap intervals. At present zero rows have
two independent human labels, so those performance estimates and human--AI
agreement are unavailable. The blinded AI comparison follows human consensus,
not the reverse. The present results establish executable coverage, workload,
and diagnostic mechanisms; they do not certify the quality of every rule.
